\documentclass[12pt,a4paper]{article}
\usepackage[utf8]{inputenc}
\usepackage[indonesian]{babel}
\usepackage{amsmath}
\usepackage{amsfonts}
\usepackage{amssymb}
\usepackage{geometry}

\geometry{margin=2.5cm}

\title{Menghitung Volume Benda Putar dengan Metode Kulit Tabung}
\author{Ahmad Fakhrul Bawani}
\date{}

\begin{document}

\maketitle

\section*{Soal}
Use the shell method to find the volumes of the solids generated by revolving the shaded region about the indicated axes:
\begin{center}
  \textbf{a.} The $x$-axis \quad \textbf{b.} The line $y = 5$ \quad \textbf{c.} The line $y = 6$ \quad \textbf{d.} The line $y = -\frac{3}{8}$
\end{center}
Daerah dibatasi oleh kurva kanan $x = \frac{2y^2}{5}$ dan kurva kiri $x = \frac{y^4}{25} - \frac{3y^2}{5}$ dari $y = 0$ hingga $y = 5$.

---

\subsection*{Rumus Dasar Tinggi Elemen}
Tinggi elemen horizontal kulit tabung ($h$) diperoleh dari pengurangan kurva kanan dikurangi kurva kiri:
\begin{equation*}
  h = \frac{2y^2}{5} - \left(\frac{y^4}{25} - \frac{3y^2}{5}\right) = y^2 - \frac{y^4}{25}
\end{equation*}
Batas integrasi untuk seluruh soal adalah dari $y = 0$ sampai $y = 5$.

---

\subsection*{a. Diputar terhadap Sumbu-$x$ (Garis $y = 0$)}
\begin{itemize}
  \item Jari-jari kulit tabung: $r = y$
  \item \textbf{Setup Integral:}
    \begin{equation*}
      V = \int_{0}^{5} 2\pi \cdot y \cdot \left(y^2 - \frac{y^4}{25}\right) \, dy = 2\pi \int_{0}^{5} \left(y^3 - \frac{y^5}{25}\right) \, dy
    \end{equation*}
  \item \textbf{Perhitungan:}
    \begin{align*}
      V &= 2\pi \left[ \frac{1}{4}y^4 - \frac{1}{150}y^6 \right]_{0}^{5} \\
      &= 2\pi \left( \frac{625}{4} - \frac{15625}{150} \right) = 2\pi \left( \frac{625}{4} - \frac{625}{6} \right) \\
      &= 2\pi \left( \frac{1875 - 1250}{12} \right) = 2\pi \left( \frac{625}{12} \right) = \frac{625\pi}{6}
    \end{align*}
\end{itemize}

\subsection*{b. Diputar terhadap Garis $y = 5$}
\begin{itemize}
  \item Jari-jari kulit tabung: $r = 5 - y$
  \item \textbf{Setup Integral:}
    \begin{equation*}
      V = \int_{0}^{5} 2\pi \cdot (5 - y) \cdot \left(y^2 - \frac{y^4}{25}\right) \, dy = 2\pi \int_{0}^{5} \left(5y^2 - \frac{y^4}{5} - y^3 + \frac{y^5}{25}\right) \, dy
    \end{equation*}
  \item \textbf{Perhitungan:}
    \begin{align*}
      V &= 2\pi \left[ \frac{5}{3}y^3 - \frac{1}{25}y^5 - \frac{1}{4}y^4 + \frac{1}{150}y^6 \right]_{0}^{5} \\
      &= 2\pi \left( \frac{625}{3} - 125 - \frac{625}{4} + \frac{625}{6} \right) \\
      &= 2\pi \left( \frac{2500 - 1500 - 1875 + 1250}{12} \right) = 2\pi \left( \frac{375}{12} \right) = \frac{125\pi}{2}
    \end{align*}
\end{itemize}

\subsection*{c. Diputar terhadap Garis $y = 6$}
\begin{itemize}
  \item Jari-jari kulit tabung: $r = 6 - y$
  \item \textbf{Setup Integral:}
    \begin{equation*}
      V = \int_{0}^{5} 2\pi \cdot (6 - y) \cdot \left(y^2 - \frac{y^4}{25}\right) \, dy = 2\pi \int_{0}^{5} \left(6y^2 - \frac{6y^4}{25} - y^3 + \frac{y^5}{25}\right) \, dy
    \end{equationitem}
  \item \textbf{Perhitungan:}
    \begin{align*}
      V &= 2\pi \left[ 2y^3 - \frac{6}{125}y^5 - \frac{1}{4}y^4 + \frac{1}{150}y^6 \right]_{0}^{5} \\
      &= 2\pi \left( 250 - 150 - \frac{625}{4} + \frac{625}{6} \right) \\
      &= 2\pi \left( 100 - \frac{625}{4} + \frac{625}{6} \right) = 2\pi \left( \frac{1200 - 1875 + 1250}{12} \right) \\
      &= 2\pi \left( \frac{575}{12} \right) = \frac{575\pi}{6}
    \end{align*}
\end{itemize}

\subsection*{d. Diputar terhadap Garis $y = -\frac{3}{8}$}
\begin{itemize}
  \item Jari-jari kulit tabung: $r = y - \left(-\frac{3}{8}\right) = y + \frac{3}{8}$
  \item \textbf{Setup Integral:}
    \begin{equation*}
      V = \int_{0}^{5} 2\pi \cdot \left(y + \frac{3}{8}\right) \cdot \left(y^2 - \frac{y^4}{25}\right) \, dy = 2\pi \int_{0}^{5} \left(y^3 - \frac{y^5}{25} + \frac{3}{8}y^2 - \frac{3y^4}{200}\right) \, dy
    \end{equation*}
  \item \textbf{Perhitungan:}
    \begin{align*}
      V &= 2\pi \left[ \frac{1}{4}y^4 - \frac{1}{150}y^6 + \frac{1}{8}y^3 - \frac{3}{1000}y^5 \right]_{0}^{5} \\
      &= 2\pi \left( \frac{625}{4} - \frac{625}{6} + \frac{125}{8} - \frac{375}{40} \right) \\
      &= 2\pi \left( \frac{625}{12} + \frac{125}{8} - \frac{75}{8} \right) = 2\pi \left( \frac{625}{12} + \frac{50}{8} \right) \\
      &= 2\pi \left( \frac{625}{12} + \frac{25}{4} \right) = 2\pi \left( \frac{625 + 75}{12} \right) = 2\pi \left( \frac{700}{12} \right) = \frac{350\pi}{3}
    \end{align*}
\end{itemize}

\subsection*{Kesimpulan Hasil Jawaban}
\begin{itemize}
  \item \textbf{a.} Volume terhadap sumbu-$x$ adalah \textbf{$\dfrac{625\pi}{6}$} satuan volume.
  \item \textbf{b.} Volume terhadap garis $y = 5$ adalah \textbf{$\dfrac{125\pi}{2}$} satuan volume.
  \item \textbf{c.} Volume terhadap garis $y = 6$ adalah \textbf{$\dfrac{575\pi}{6}$} satuan volume.
  \item \textbf{d.} Volume terhadap garis $y = -\frac{3}{8}$ adalah \textbf{$\dfrac{350\pi}{3}$} satuan volume.
\end{itemize}

\end{document}